\chapter{Benchmark results}\label{ch:results}

Each section gives the setup, the result, and the reading.
Single particles are deterministic, and each configuration was run once per code: ``one instrument'' means a single
code path measured the number.

\section{Case A: uniform field}\label{sec:caseA}
\textbf{Setup.} $B_z = \SI{1}{T}$, a \SI{1}{MeV} proton, one turn, 90 to 1440 steps per turn.
LIONS Boris, OPAL LF2 and RK4, RF-Track adaptive integrators.

\textbf{Result} (\cref{fig:caseA}).
LIONS converges at order \AlionsRadOrder{} in radius, with \AlionsRadFourteen{} at 1440 steps.
Its fitted radius equals the Boris polygon radius of \cref{eq:boris-radius} within \AlionsPolygon, and $|\vect p|$ is exact.
OPAL LF2 keeps $|\vect p|$ within \ALFTwoP; its radius is exact because its positions are chord midpoints
(\cref{sec:boris-uniform}), so only the phase discriminates, and its phase error matches the Boris closed form.
OPAL RK4 loses $|\vect p|$ by \ARkFourPThreeSixty{} per turn at 360 steps, against \ARkFourPred{} from
\cref{eq:rk4-decay}.
Its fitted radius error is \ARkFourRadThreeSixty, half of that.
This is expected if the radius decays linearly over the turn and the fit returns the mean, but it was derived after the
run.

\textbf{Reading.}
LIONS Boris behaves exactly as the analysis predicts, to rounding.
The registered criterion for RK4 (fourth order in radius) was wrong; RK4 radius and $|\vect p|$ errors are fifth order.

\begin{figure}[t]
\centering
\includegraphics[width=0.95\textwidth]{case_a}
\caption{Case A. Left: relative radius error against steps per turn, with the Boris prediction $(\omega\Delta t)^2/8$
and the RK4 amplitude decay. Right: relative error of the angle swept per step, with the Boris closed form
$1-2\arctan(\omega\Delta t/2)/\omega\Delta t\approx(\omega\Delta t)^2/12$. LIONS and OPAL LF2 overlap in phase.}
\label{fig:caseA}
\end{figure}

\section{Case B0: smooth radial gradient}\label{sec:caseB0}
\textbf{Setup.} $B_z = B_0(1-\kappa(r-r_0))$, $B_0 = \SI{-1}{T}$, $\kappa = \SI{0.5}{m^{-1}}$, $r_0 = \SI{0.15}{m}$, 1 to \SI{5}{MeV}
protons, on a bilinear map of \SI{5}{mm} by \SI{1}{\degree}.
$B_z$ is linear in $r$ and constant in $\theta$, so bilinear interpolation is exact for it.
Exact references from \cref{sec:weak}.

\textbf{Result} (\cref{fig:caseB0}).
LIONS converges at order \BzeroLionsOrderMin{} or better, with closed-orbit errors up to \BzeroLionsFourteenMax{} at
1440 steps.
Its $\Qz$ is low by up to \BzeroLionsQzErrMax.
This is consistent with the chord error of the interpolated $\partial_xB$, $\partial_yB$ (the horizontal field
follows the chord of $(\cos\theta,\sin\theta)$ between nodes), which shifts $\Qz$ by about $10^{-5}$ relative (estimate in \nolinkurl{benchmarks/DESIGN.md}, section 12).
RF-Track reaches the tolerance level.
OPAL LF2 and RK4 share an energy-dependent orbit offset, up to \BzeroOpalOffset{} relative (RK4), that is not explained; single-precision field
storage was ruled out from the source.
The Jacobian tunes and the Fourier tunes of LIONS agree within \BzeroFFTAgree.

\textbf{Reading.}
On a field without kinks in the plane, LIONS is a clean second-order code.
The orbit-error tolerance of $10^{-6}$ in the first draft is out of reach of Boris at 1440 steps, whose radius error in a
uniform field is $(\omega\Delta t)^2/8 = \AlionsRadFourteen$.
It was dropped after an exploratory LIONS run at \SI{1}{MeV} showed this, and before the registered B0 runs.

\begin{figure}[t]
\centering
\includegraphics[width=0.6\textwidth]{case_b0}
\caption{Case B0: closed-orbit radius error against the exact value, range over 1 to \SI{5}{MeV} (band) and maximum (line).}
\label{fig:caseB0}
\end{figure}

\section{Case B: 4-sector field on a bilinear map}\label{sec:caseB}
\textbf{Setup.} $B_z = B_0(1+0.3\cos4\theta)$, $B_0 = \SI{-1}{T}$, same map grid.
Report only: no pass/fail.

\textbf{Result} (\cref{fig:caseB}).
The tunes scatter with the step count, by up to \BQrSpreadMax{} in $\Qr$ over the map codes (LIONS:
\BLionsQrSpread), with no trend.
$|\det M - 1|$ ranges from \BDetMin{} to \BDetMax{} for the map codes, also without a trend.
On the closed-form field without a map, RF-Track at the same tune scatters by \BsmoothQrSpread{} in $\Qr$.
Halving the azimuthal map step does not reduce the LIONS scatter (ratio \BHalfDegQrRatioMin{} to
\BHalfDegQrRatioMax).

\textbf{Discriminating runs} (case B2, registered before running).
A radial gradient moves $\Qr$ away from 1.
Per energy, $\det(M_r - I)$ grows by \BTwoCondMin{} to \BTwoCondMax{} times (from the LIONS $\Qr$ at 1440 steps), so the
closed-orbit sensitivity $1/(2\sin\pi\Qr)$ falls by \BTwoSensMin{} to \BTwoSensMax{} (\cref{sec:conditioning}).
The orbit spread across step counts fell by \BTwoSpreadLionsMin{} to \BTwoSpreadLionsMax{} (LIONS), \BTwoSpreadLFMin{} to
\BTwoSpreadLFMax{} (OPAL LF2) and \BTwoSpreadRKMin{} to \BTwoSpreadRKMax{} (OPAL RK4).
On the closed-form field RF-Track stayed clean at the near-integer tune.

\textbf{Reading, corrected.}
The registered prediction expected the spread to fall by the full $\det(M_r-I)$ ratio, and on that basis
\nolinkurl{benchmarks/DESIGN.md} (section 15) rejected H4.
The prediction used the wrong scaling: the sensitivity goes as $1/\sqrt{\det(M_r-I)}$.
With the correct scaling, the OPAL spreads fall by amounts of the predicted size, within a factor of a few either way,
while the LIONS spread falls as predicted at one energy of five and much less at the others.
The clean RF-Track run shows that the map is the source of the errors; it says nothing about amplification, because the
errors it would amplify are tiny.
So the map is the source, and near-integer amplification (H4) is not excluded: it may add to the scatter at low
energy.

\begin{figure}[t]
\centering
\includegraphics[width=0.6\textwidth]{case_b_scatter}
\caption{Case B at \SI{1}{MeV}: fractional $\Qr$ against steps per turn.
The map codes scatter at the $10^{-3}$ level; RF-Track on the closed-form field (no map) is flat.}
\label{fig:caseB}
\end{figure}

\section{Case D2: a long-term invariant}\label{sec:caseD}
\textbf{Setup.}
In the B0 field the canonical angular momentum is an exact invariant of planar motion.
In LIONS units,
\begin{equation}
L = x p_y - y p_x + \frac{q}{m_0c}\,B_0\left(\frac{(1+\kappa r_0)\,r^2}{2} - \frac{\kappa\,r^3}{3}\right).
\end{equation}
One particle at \SI{1}{MeV}, \SI{1}{mm} off the closed orbit, \DTurns{} turns at 360 steps per turn (the largest step).
A secular change below $10^{-9}$ counts as no drift.

\textbf{Result} (\cref{tab:caseD}).
LIONS Boris and OPAL LF2 show no drift and a bounded oscillation of $L$; $|\vect p|$ is conserved to rounding.
RK4 loses $|\vect p|$ as \cref{eq:rk4-decay} predicts over \DTurns{} turns (\DRkFourPred), and $L$ drifts with it.
RF-Track drifts at the level of its tolerance accumulated over the run.

\textbf{Reading.}
On this test the Boris step behaves like a symplectic integrator: bounded oscillation, no secular drift.
That is consistent with symplectic behaviour; it is not a proof.
One field, one amplitude, planar motion.
The LIONS oscillation of $L$ is \DOscRatio{} times that of LF2; this was not analysed.

\begin{table}[h]
\centering\small
\caption{Case D2, \DTurns{} turns at 360 steps per turn.}
\label{tab:caseD}
\input{gen/tab_d2.tex}
\end{table}

\section{Case G1: manufactured solutions}\label{sec:caseG1}
\textbf{Setup.}
FROM\_ORBIT maps from the published 3-sector (H$_2^+$, three orbits from 19 to \SI{130}{MeV}) and 10-sector examples
(protons, four orbits from 932 to \SI{1630}{MeV}) of FROM\_ORBIT \cite{planche2019,planche2023}, shown in
\cref{fig:field-maps}.
The exact field is evaluated by FROM\_ORBIT's own root finding on a \thirddeg{} master grid; coarser maps are
subsets of its nodes.
\begin{itemize}
\item Step 1: RF-Track on a quintic spline of the master measures the precision of the references.
On the 3-sector map the prescribed orbits are reproduced to \GThreeRefOrbit{} and the FROM\_ORBIT geometric tunes are
good to \GThreeRefQr{} ($\Qr$) and \GThreeRefQz{} ($\Qz$); on the 10-sector map, to \GTenRefOrbit{} in orbit radius and
\GTenRefQr{} in the tunes.
The RF-Track tunes become the tune references.
\item Step 2: LIONS on bilinear maps at \thirddeg{}, \SI{1}{\degree} and \SI{4}{\degree}.
\item Step 3: the same Boris step (numpy harness, validated against LIONS) on bicubic and quintic splines of the same nodes.
\end{itemize}

\textbf{Result} (\cref{fig:G1}).
At \SI{4}{\degree} LIONS could not find the closed orbit in \GThreeFourDegFailed{} (3-sector) and \GTenFourDegFailed{}
(10-sector) configurations.
At \thirddeg{} and 1440 steps, the worst errors are:

\begin{center}\small
\begin{tabular}{lcccc}
\toprule
 & \multicolumn{2}{c}{3-sector} & \multicolumn{2}{c}{10-sector}\\
 & bilinear (LIONS) & spline (Boris) & bilinear (LIONS) & spline (Boris)\\
\midrule
$|\Delta\Qr|$ & \GThreeBilQrMax & \GThreeSplQrMax & \GTenBilQrMax & \GTenSplQrMax\\
$|\det M-1|$ & \GThreeBilDetMax & \GThreeSplDetMax & \GTenBilDetMax & \GTenSplDetMax\\
$|\Delta r_{eo}|/r$ & \GThreeBilReoMax & \GThreeSplReoMax & \GTenBilReoMax & \GTenSplReoMax\\
\bottomrule
\end{tabular}
\end{center}

With splines the errors converge at second order in orbit, turn time and tunes (closed-orbit error ratio per halving of
the step \GThreeSplOrbitRatioMin{} to \GThreeSplOrbitRatioMax{} on the 3-sector map, \GTenSplOrbitRatioMin{} to
\GTenSplOrbitRatioMax{} on the 10-sector map).
With the bilinear field the tune errors do not converge on the 3-sector map, and converge at most at first order on the
10-sector map (worst $|\Delta\Qr|$ \GTenBilQrMaxThreeSixty{} at 360 steps, \GTenBilQrMax{} at 1440); $\det M - 1$ stays
orders of magnitude above the spline values.
At 360 steps the closed-orbit errors of both are close, so there the integrator dominates.
At 1440 steps the bilinear orbit error is larger than the spline one on the 10-sector map.
On the 10-sector map at \SI{1}{\degree} (18 nodes per \SI{18}{\degree} harmonic) even the splines stop converging
cleanly.

\textbf{Reading.}
Only the interpolation changed between steps 2 and 3, so the tune excess and the $\det$ excess on these maps come with the
bilinear interpolation; \cref{tab:evidence} separates the field model from the integration across its kinks.
With splines, the isochronism error (turn time) at 1440 steps is set by the Boris step: at most \GThreeSplTurnTimeMax{}
(3-sector) and \GTenSplTurnTimeMax{} (10-sector); on the bilinear maps it is larger.
Over 500 turns at harmonic number $h_{\rm rf}$, a turn-time error of $10^{-6}$ would be an RF phase slip of order
$2\pi h_{\rm rf}\times500\times10^{-6}$~rad; this is an estimate, not simulated.

\begin{figure}[t]
\centering
\includegraphics[width=0.9\textwidth]{g1}
\caption{Case G1 at \thirddeg{}: $|\Delta\Qr|$ and $|\det M-1|$ against steps per turn.
LIONS on the bilinear map (blue) against the same Boris step on bicubic and quintic splines of the same nodes.
Bands: range over the orbits; lines: maximum.}
\label{fig:G1}
\end{figure}

\section{Case G2: PSI Ring against FIXPO}\label{sec:caseG2}
\textbf{Setup.}
The PSI Ring \SI{590}{MeV} map (8-fold, \thirddeg{}, $r$ from 1.9 to \SI{4.7}{m}) and the FIXPO equilibrium orbits
and tunes from 72 to \SI{530}{MeV}.
LIONS on the native-resolution map (1081 columns), the Boris harness on a quintic spline, RF-Track on the same spline,
and the OPAL closed-orbit finder.
All at 1440 steps per turn (RF-Track: 2880 samples), no trim coil.
At \SI{512}{MeV} OPAL returned a non-solution without an error; that energy is excluded for every code in the table.

\textbf{Result} (\cref{fig:G2}, \cref{tab:G2}).
Closed-orbit radii of every code agree with FIXPO at the level of the print resolution of the reference file
(\SI{0.1}{mm}), so the radius comparison cannot separate the codes.
The smooth-field instruments are closest to FIXPO in $\Qr$ (median \GPsiSplineQrMed{} and \GPsiRftrackQrMed), then OPAL
(\GPsiOpalQrMed), then LIONS (\GPsiLionsQrMed).
On the bilinear map the radial tune is defined only roughly: the case H harness with the same node data (F0, Boris,
$\operatorname{atan2}$ azimuth, its own Newton solution) agrees with LIONS at some energies and differs several-fold at
others, giving a median of \HFixpoBilQrMed{} against FIXPO (\cref{sec:caseH}).
LIONS has $|\det M - 1|$ of \GPsiLionsDetMed{} (median), up to \GPsiLionsDetMax.
All these tunes come from the one-turn Jacobian.
Tracked over many turns, LIONS on the same field is within \GThreeFiveFixMax{} of FIXPO at \SI{5}{mm} (\cref{sec:caseG3}), so the LIONS line of \cref{tab:G2} does not measure the betatron frequency of a beam.

\begin{table}[h]
\centering\small
\caption{Case G2: absolute errors against FIXPO, median / maximum over the interior energies of the step-3 grid
(92 to \SI{492}{MeV}, without the 386 to \SI{430}{MeV} band and \SI{512}{MeV}).}
\label{tab:G2}
\input{gen/tab_g2.tex}
\end{table}

\begin{figure}[t]
\centering
\includegraphics[width=0.8\textwidth]{g2_tunes}
\caption{Case G2: tune errors against FIXPO over the energy range.
Grey band: $\Qr$ near 1.5, where the trace formula is ill-conditioned.}
\label{fig:G2}
\end{figure}

\textbf{Discriminating runs} (steps 2b and 2c, \cref{tab:G2det}).
The numpy Boris on the bilinear field reproduces LIONS, so the effect is not a LIONS defect (H2 rejected).
The determinant barely moves with the Jacobian step, which argues against a finite-difference artefact (H3) without
excluding it (\cref{ch:orbits}, \cref{tab:evidence}).
The effect is in the radial block only; the vertical block is clean.
RF-Track (adaptive) on the same bilinear field gives a $|\det M_r - 1|$ \GTwoCRatioMin{} to \GTwoCRatioMax{} times
smaller; the registered reading threshold ($10^{-3}$) was met at one energy of three.
Since the exact flow preserves area on any continuous field (\cref{ch:integrators}), the $\det$ excess must come from
the discretisation.
This is consistent with H1b, fixed-step integration across gradient jumps; case H shows that fixed-step RK4 has the
same excess (\cref{sec:caseH}).

\begin{table}[h]
\centering\small
\caption{Case G2: $\det M_r - 1$ of the radial block at three energies, Jacobian step \SI{1e-6}{m}.}
\label{tab:G2det}
\input{gen/tab_g2_det.tex}
\end{table}

\paragraph{A criterion that could not fail.}
The registered G2 criterion was $|L-F|\le|O-F|+|L-S|$ per metric and energy, with $L$ LIONS, $F$ FIXPO, $O$ OPAL and $S$
the spline tracker.
By the triangle inequality $|L-F|\le|L-S|+|S-F|$, so whenever the spline tracker is at least as close to FIXPO as OPAL,
the criterion holds whatever LIONS does.
It passed for $\Qr$ while LIONS was \GPsiLionsOverOpal{} times further from FIXPO than OPAL (median).
The informative comparison is \cref{tab:G2}.
The criterion was proposed by this analysis, and the record says so.

\section{Case G3: OPALX and the step phase}\label{sec:caseG3}
\textbf{Question} (registered before running).
Does OPALX 1.0.0 \cite{opalx2026} reproduce OPAL 2024.1 and FIXPO on the PSI map, and does a third, independent tracker on the same bilinear field show the $\det M_r$ excess?
OPALX interpolates $B_z$ bilinearly on the same nodes.
Its closed-orbit command writes the Jacobian; no trim coil, 1440 steps, the G2 step 3 energies.
The installation reproduces its own regression test on this map within its tolerances.

\textbf{Result} (\cref{tab:G3,tab:G3det}).
No run failed.
The excess is not specific to LIONS: OPALX Boris has a median $|\det M_r-1|$ of \GThreeBorisDet{}, RK4 \GThreeBorisOverRk{} times less, and the fixed-step eighth-order DOP853 \GThreeDopOverRk{} times more than RK4.
No integrator shows its order: the values are not monotonic in the step count and they change sign.
The OPALX manual reports the same behaviour on this map at 46080 steps per turn, with a radial modulus off unity by $10^{-4}$ to $10^{-3}$, and names the piecewise interpolation as an unconfirmed candidate \cite{opalx2026}.
OPALX Boris $\det M_r$ moves at most \GThreeFdVar\,\% with the derivative step from $10^{-7}$ to $10^{-5}$.

\begin{table}[h]
\centering\small
\caption{Case G3: OPALX against FIXPO, median / maximum over the energies of \cref{tab:G2}, Jacobian tunes, 1440 steps.
$|\operatorname{tr}M_r/2|>1$ at several energies, where the tune is clipped to 0.5; the maxima include these.}
\label{tab:G3}
\input{gen/tab_g3.tex}
\end{table}

\begin{table}[h]
\centering\small
\caption{Case G3: OPALX $\det M_r - 1$ against steps per turn.}
\label{tab:G3det}
\input{gen/tab_g3_det.tex}
\end{table}

One registered prediction failed.
OPALX minus OPAL 2024.1 in $\Qr$ has a median of \GThreeOpalQrMin{} to \GThreeOpalQrMax{}, where below $10^{-3}$ was predicted, because the premise of a shared field model does not hold for tunes.
From its source, the OPAL finder integrates the linearised equations with the interpolated $\partial_rB$ and $\partial_\theta B$ channels, so it sees a continuous gradient.
OPALX differentiates the tracked map, which sees the gradient jumps of bilinear $B_z$.
This explanation is read from the sources, not tested.

\paragraph{The same field in two codes.}
LIONS and OPALX Boris have the same magnitude of $\det M_r-1$ but not the same value: the correlation over the 24 energies is \GThreeCorr.
OPALX also subdivides steps at sector edges.
So $\det M_r$ is a property of the discrete map and of where its steps fall, not a stable property of the field.

\paragraph{Tracked tunes.}
OPALX spectral tunes (a ray \SI{5}{mm} off, RK4, 100 turns) lie within \GThreeSpecQrMax{} of FIXPO in $\Qr$ at \GThreeSpecN{} of \GThreeSpecN{} energies (bin \num{2.5e-3}), the half-integer band included.
A registered run then measured LIONS itself: launch at amplitude $a$ in $r$ and $z$ from the code's own closed orbit, 100 turns, tune from the turn-by-turn signal (\cref{tab:G3amp}, \cref{fig:G3}).
LIONS is within \GThreeFiveSplMax{} of the spline tracker at \SI{5}{mm}, \GThreeOneSplMax{} at \SI{1}{mm} and \GThreeTenthSplMax{} at \SI{0.1}{mm}.
At \SI{0.01}{mm} it is more than 0.1 off at \GThreeSmallOff{} of \GThreeSmallN{} energies.
The Hermite option stays close at every amplitude (median against FIXPO \GThreeHermQrMed{} in $\Qr$ and \GThreeHermQzMed{} in $\Qz$; spline \GThreeSplQzMed{} and default \GThreeLionsQzMed{} in $\Qz$).

\begin{table}[h]
\centering\small
\caption{Case G3 step 2: tracked $|\Qr-\Qr^{\rm FIXPO}|$, median / maximum over 22 energies outside the band, and the ratio of the rms radial amplitude over the last and first ten of 100 turns.
Baseline of that ratio on the smooth models: \GThreeArBaseMin{} to \GThreeArBaseMax{} (beating at low tune).}
\label{tab:G3amp}
\input{gen/tab_g3_amp.tex}
\end{table}

\begin{figure}[t]
\centering
\includegraphics[width=\textwidth]{g3_phase}
\caption{Left: radial tune error of LIONS on the PSI map, from the one-turn Jacobian and from tracking.
Grey band: $\Qr$ near 1.5.
Right: distance from the closed orbit at \SI{292}{MeV}, in units of the launch amplitude, at the nominal step and on a commensurate turn.}
\label{fig:G3}
\end{figure}

\paragraph{The step phase.}
A \SI{0.01}{mm} oscillation does not grow as $\det M_r$ implies: its amplitude ratio is at most \GThreeSmallArMax, against up to \GThreeSmallPred{} from $(\det M_r)^{45}$.
The registered reading of that failure was that the Jacobian was suspect again, so the Jacobian was tested directly (step 2b, 30 turns, \SIrange{1e-7}{1e-4}{m}).
$M$ predicts the first return within \GThreeOneTurnMin{} to \GThreeOneTurnMax{} at $a=\SI{1e-7}{m}$.
At the second return the prediction fails by order 1, which also missed the registered prediction.
From then on the distance to the closed orbit no longer depends on $a$: at \SI{152}{MeV} and turn 5 it is \GThreeJitterA{}~m from $a=\SI{1e-7}{m}$ and \GThreeJitterB{}~m from $a=\SI{1e-6}{m}$.

Hypothesis S: the turn time is not an integer number of steps, so each turn starts at another step phase and is another map.
At the nominal step these energies miss a commensurate turn by \GThreeSlipMin{} to \GThreeSlipMax{} of a step per turn (turn time between the first two crossings, G2 step 2).
Test (step 2c): adjust the step until one turn is exactly 1440 steps, and repeat.
At turn 2 the deviation falls to \GThreeComTwoMin{} to \GThreeComTwoMax, and the growth no longer depends on $a$ and tracks $M^n$ (\cref{fig:G3}, right).
At \SI{292}{MeV} a launch grows by \GThreeComGrow{} in 20 turns, against \GThreeComGrowPred{} from $M^n$.
This refutes the rival that the second-return failure does not depend on the step phase.
It misses the registered threshold ($10^{-2}$ through turn 10): the deviation at turn 10 is \GThreeComTenMin{} to \GThreeComTenMax, not explained.

\begin{reading}
This reading is a proposal.
On a bilinear field, the one-turn map depends on the step phase at the section.
The Newton solution and $M$ of cases G2, H and G3 are correct for one phase.
Over many turns the orbit moves among these maps.
Then:
(1) tracked tunes are right, with a median of \GThreeTenthFixMed{} against FIXPO at \SI{0.1}{mm};
(2) below about \SI{0.1}{mm} the betatron signal is lost in a jitter of tens of micrometres;
(3) amplitudes are not reliable: the radial amplitude at \num{0.1} and \SI{1}{mm} changes by \GThreeArMin{} to \GThreeArMax{} times in 100 turns on a coasting orbit;
(4) a deliberately commensurate step makes the growth exponential;
(5) the Hermite field removes all of this in the runs above.
That each phase has its own fixed point is inferred, not measured.
Beam size and emittance over longer runs, on a bunch and with acceleration have not been measured.
\end{reading}
